\documentclass[11pt,a4paper]{article}

% Packages
\usepackage[utf8]{inputenc}
\usepackage[T1]{fontenc}
\usepackage[margin=0.75in]{geometry}
\usepackage{enumitem}
\usepackage{hyperref}
\usepackage{titlesec}
\usepackage{xcolor}

% Formatting
\pagestyle{empty}
\setlength{\parindent}{0pt}
\setlist[itemize]{leftmargin=*, noitemsep, topsep=3pt}

% Section formatting
\titleformat{\section}{\Large\bfseries}{}{0em}{}[\titlerule]
\titlespacing*{\section}{0pt}{12pt}{6pt}

\titleformat{\subsection}[runin]{\bfseries}{}{0em}{}
\titlespacing*{\subsection}{0pt}{8pt}{0pt}

% Colors
\definecolor{linkcolor}{RGB}{0,102,204}
\hypersetup{
    colorlinks=true,
    urlcolor=linkcolor
}

\begin{document}

% Header
\begin{center}
    {\LARGE \textbf{ERIC MAYEUX}}\\
    \vspace{4pt}
    {\large Gestionnaire de produit principal -- Opérations produit | Lead Product Operations Manager}\\
    \vspace{4pt}
    Montréal, QC, Canada\\
    \href{mailto:mayeux.e@gmail.com}{mayeux.e@gmail.com} | 438-884-7645
\end{center}

\vspace{8pt}

\section*{Profil Professionnel}
\noindent
\textbf{Gestionnaire de produit principal} avec 10+ ans d'expérience en \textbf{transformation produit} et \textbf{coaching de Propriétaires de produit} en environnement agile. Expert en \textbf{modèle opérationnel orienté produit}, priorisation guidée par les données et \textbf{amélioration continue}. Maîtrise des pratiques de Product Ops : Productboard, Jira, Confluence, \textbf{tableaux de bord KPI}. Certifié \textbf{PMP} (en cours), CSM. Bilinguisme français/anglais.

\section*{Compétences Techniques}

\textbf{Product Operations \& Transformation:} Modèle opérationnel orienté produit, coaching PO, discovery continue, priorisation (RICE, WSJF), standards de livraison, communauté de pratique

\textbf{Leadership sans autorité:} Mentorat de Propriétaires de produit, shadowing, rétroaction en temps réel, amélioration des pratiques d'équipe, formation

\textbf{Facilitation \& Communication:} Animation d'ateliers, storytelling, présentation aux parties prenantes, alignment multi-équipes, bilinguisme français/anglais

\textbf{Outils:} Jira, Jira Product Discovery, Confluence, Datadog, Splunk, tableaux de bord KPI, Productboard

\textbf{Méthodologies:} Agile, Scrum, SAFe, Kanban, développement itératif, amélioration continue, PMP

\textbf{Intelligence Artificielle:} Agents IA, Playwright MCP, Crew AI, automatisation, RPA

\textbf{Architecture \& Cloud:} AWS Cloud Practitioner, CI/CD, Architectures Distribuées

\vspace{6pt}

\section*{Réalisations}
\begin{itemize}
    \item \textbf{Coaching et mentorat} de 5+ professionnels TI en contexte agile SAFe à la BNC
    \item Mise en place d'un \textbf{modèle opérationnel} qualité : standards, processus, métriques, tableaux de bord
    \item Définition des \textbf{pratiques de priorisation} et découverte continue pour les équipes de livraison
    \item \textbf{Storytelling} : présentation des KPIs DORA, DevEx et qualité aux parties prenantes et à la direction
    \item Organisation de hackathons et POC : \textbf{amélioration continue} des pratiques de développement
    \item Déploiement d'\textbf{agents IA} pour l'automatisation des tâches de gestion et de reporting
\end{itemize}

\section*{Expérience Professionnelle}

\subsection*{Scrum Master -- Gestion de projet | Project Manager} \hfill \textit{Novembre 2025 -- Présent}\\
\textbf{Banque Nationale du Canada (BNC)} -- Montréal, QC\\
Pilote la \textbf{transformation produit} et coache les équipes vers un \textbf{modèle orienté produit}.
\begin{itemize}
    \item Observe et coache les équipes IT ; fournit des \textbf{conseils pratiques et rétroaction en temps réel}
    \item Forme sur les compétences produit : \textbf{discovery}, priorisation, engagement des parties prenantes, storytelling
    \item Définit et renforce les \textbf{pratiques du modèle opérationnel} produit, de la discovery à la livraison
    \item Opérationnalise des \textbf{tableaux de bord KPI} (DORA, DevEx, qualité) pour les escouades
    \item Déploie des outils d'\textbf{amélioration continue} et leçons apprises à l'échelle des équipes
    \item Exploite des \textbf{agents IA} pour l'automatisation des tâches administratives et de gestion
\end{itemize}

\subsection*{Intégrateur Sénior Qualité | Senior Quality Integrator} \hfill \textit{Octobre 2023 -- Novembre 2025}\\
\textbf{Banque Nationale du Canada (BNC)} -- Montréal, QC\\
\begin{itemize}
    \item \textbf{Coaching} et mentorat des QA ; recrutement, onboarding, formation continue
    \item Définit le \textbf{modèle opérationnel} qualité pan-train : standards, métriques, gouvernance
    \item \textbf{Évalue la maturité} des équipes et met en place des plans d'amélioration ciblés
\end{itemize}

\subsection*{Intégrateur Qualité -- SDET | Quality Integrator -- Software Development Engineer in Test} \hfill \textit{Janvier 2022 -- Octobre 2023}\\
\textbf{Banque Nationale du Canada (BNC)} via Groupe Astek -- Montréal, QC\\
\begin{itemize}
    \item Mise en place de \textbf{cadres de priorisation} qualité et standards de préparation
    \item Stack : Selenium, AppliTools, RestAssured, JMeter, Gatling, K6
\end{itemize}

\subsection*{QA Automaticien | QA Automation Engineer} \hfill \textit{Juin 2021 -- Décembre 2021}\\
\textbf{AODocs} -- Paris, France\\
\begin{itemize}
    \item POC et \textbf{amélioration continue} des pratiques de test automatisé
\end{itemize}

\subsection*{Automaticien CI/CD | DevOps Automation Engineer} \hfill \textit{Janvier 2020 -- Juin 2021}\\
\textbf{ENGIE DIGITAL} via Groupe Hénix -- Paris, France\\
\begin{itemize}
    \item Déploiement CI/CD ; \textbf{administration Jira}, KPIs, pratiques DevOps
\end{itemize}

\subsection*{Product Owner} \hfill \textit{Juin 2019 -- Septembre 2019}\\
\textbf{Hénix} -- Montrouge, France\\
\begin{itemize}
    \item \textbf{Gestion du backlog}, priorisation, certification Jira ACP-600
\end{itemize}

\subsection*{Développeur Logiciel | Software Developer} \hfill \textit{Mai 2017 -- Mai 2019}\\
\textbf{MEDIAPOST} via Groupe Hénix -- Paris, France\\
\begin{itemize}
    \item Développement back office, reporting ; expérience du cycle de livraison de A à Z
\end{itemize}

\subsection*{Consultant QA \& Automatisation | QA Automation Consultant} \hfill \textit{Avril 2016 -- Avril 2017}\\
\textbf{ENGIE IT} via Groupe Hénix -- Saint-Ouen, France\\
\begin{itemize}
    \item Coordination opérationnelle, administration HP ALM/UFT
\end{itemize}

\section*{Certifications \& Formations}

\textbf{PMP Certification} (en cours) \hfill \textit{2026}

Project Management Professional -- Project Management Institute (PMI)

\textbf{AWS Cloud Practitioner} \hfill \textit{2026}

Amazon Web Services (AWS)

\textbf{Certified Scrum Master (CSM)} \hfill \textit{2024}

Scrum Alliance

\textbf{Jira ACP-600 -- Project Administration in Jira Server} \hfill \textit{2019}

Atlassian Certified Professional

\textbf{Cursus Automaticien et DevOps} \hfill \textit{2019}\\
\textit{AFCEPF -- Montrouge, France}

\textbf{Cursus Architecture logiciel \& Analyste informaticien - Certifié Master II} \hfill \textit{2017}\\
\textit{AFCEPF-HENIX -- Malakoff, France}

\section*{Formation Académique}

\textbf{Licence en Gestion (Bachelor's Degree in Management)} \hfill \textit{2011}\\
École Supérieure de Gestion (E.S.G) -- Paris, France

\section*{Compétences Linguistiques}

\textbf{Français:} Langue maternelle (Native)\\
\textbf{Anglais:} Niveau Avancé (Advanced / Professional Working Proficiency)

\end{document}
